%% =================================================================
%% Wei, Mei, Zhao, Xiao, Sun, Zhang, Guo.
%% "Nondestructively identifying the mechanical behavior of soft
%%  tissues using surface deformation with an explicit inverse
%%  approach."
%% Appl. Math. Modelling 134 (2024) 126--147.
%%
%% Reconstruction of page 21 (printed p. 146) as study notes.
%% Generated by: reproduce-basic-paper-tex → reproduce-multiple-pages-basic-tex [17, 22]
%%
%% Structure: Fig. 27 (varying inclusion shapes) + end-matter
%% sections (CRediT, competing interest, data availability,
%% acknowledgments) + References [1]-[4].
%% No equations. This is the first time the Wei reconstruction
%% has hit a References-as-content page.
%% =================================================================

\documentclass[11pt]{article}

\usepackage[margin=1in]{geometry}
\usepackage{amsmath, amssymb}
\usepackage{mathrsfs}
\usepackage{bm}
\usepackage{xcolor}
\usepackage{soul}
\definecolor{hlyellow}{RGB}{255,255,170}
\definecolor{hlcyan}{RGB}{170,255,255}
\definecolor{hlgreen}{RGB}{170,255,170}
\definecolor{hlpink}{RGB}{255,170,170}
\definecolor{hlorange}{RGB}{255,221,170}
\usepackage{fancyhdr}

\pagestyle{fancy}
\fancyhf{}
\lhead{\small Y.~Mei et al.}
\rhead{\small Applied Mathematical Modelling 134 (2024) 126--147}
\cfoot{\small 146 $\to$ \arabic{page}}
\renewcommand{\headrulewidth}{0.4pt}

\setcounter{page}{1}

\begin{document}

% ---------- Figure 27 placeholder ----------
\noindent
\begin{center}
\fbox{\parbox{0.92\textwidth}{\centering\vspace{1em}
\textbf{[Figure 27 --- placeholder]}\\[0.8em]
Six-panel comparison of \textbf{two different inclusion shapes}
reconstructed with 3 displacement datasets (Sample 1, no noise).
The two rows show different target shapes; the columns show
target, explicit method result, and implicit method result.\\[0.4em]
\textbf{Top row (irregular/elongated inclusion)}:
(a) target: red elongated/irregular shape, SM $\in [1.00, 5.00]$\,MPa;
(b) Explicit, 3 fields: SM $\in [1.00, 5.00]$\,MPa --- the
reconstructed shape faithfully preserves the irregular target
geometry;
(c) Implicit, 3 fields: SM $\in [1.00, 4.98]$\,MPa --- the shape
is \emph{distorted}, rendered as a green/yellow blob without
the sharp edges of the target, and fails to capture the elongated
profile.\\[0.4em]
\textbf{Bottom row (square inclusion)}:
(d) target: red square shape, SM $\in [1.00, 5.00]$\,MPa;
(e) Explicit, 3 fields: SM $\in [1.00, 5.00]$\,MPa --- the
reconstructed shape is a rounded-square, closely matching the
target;
(f) Implicit, 3 fields: SM $\in [1.00, 4.71]$\,MPa --- the shape
is rendered as a green amorphous blob, significantly distorting
the square target.\\[0.3em]
\textbf{Takeaway}: this figure is the qualitative case for the
explicit inverse method's main selling point on non-circular
inclusion geometries. The explicit method's B-spline-based shape
parameterization allows it to faithfully represent elongated or
square shapes; the implicit method's pixel-based optimization
reverts to blob-like reconstructions regardless of the target
shape. This is a shape-fidelity argument, not a shear-modulus
accuracy argument --- and it's where the explicit method's prior
knowledge (topology + rough shape) pays off most visibly.
\vspace{1em}}}
\end{center}
\vspace{0.3em}
\begin{center}
\textbf{Fig. 27.} The target and reconstructed results for
varying shapes of inclusion with 3 displacement datasets
corresponding to Sample 1 (Without noise, Unit: MPa).
\end{center}

\vspace{1em}

% ---------- CRediT authorship contribution statement ----------
\noindent\textbf{CRediT authorship contribution statement}

\vspace{0.3em}

\noindent
\textbf{Yue Mei:} Writing -- review \& editing, Writing --
original draft, Methodology, Formal analysis, Conceptualization.
\textbf{Dongmei Zhao:} Writing -- review \& editing, Writing --
original draft, Software, Formal analysis, Data curation.
\textbf{Changjiang Xiao:} Writing -- original draft,
Investigation, Formal analysis, Data curation.
\textbf{Zhi Sun:} Methodology, Formal analysis, Data curation.
\textbf{Weisheng Zhang:} Writing -- review \& editing, Writing
-- original draft, Supervision, Funding acquisition.
\textbf{Xu Guo:} Writing -- review \& editing, Writing --
original draft, Supervision, Funding acquisition.

\vspace{0.5em}

% ---------- Declaration of competing interest ----------
\noindent\textbf{Declaration of competing interest}

\vspace{0.3em}

\noindent
The authors declare that they have no known competing financial
interests or personal relationships that could have appeared to
influence the work reported in this paper.

\vspace{0.5em}

% ---------- Data availability ----------
\noindent\textbf{Data availability}

\vspace{0.3em}

\noindent
\pdfcomment[icon=Comment,color=green]{User input: that they haven't provided publicly available data in the Wei paper\\\\Claude comments: SUPPORTS the claim. "Made available on request" is explicitly NOT "publicly available." Publicly available would mean deposited in an open repository (GitHub, Zenodo, institutional archive) with a URL anyone can access. "On request" means you must contact the authors, who may or may not respond, may impose conditions, and provide no guarantee of timely access. This is a common practice in computational mechanics but falls short of modern open-science standards.\\\\Additionally, the paper provides NO code/software availability statement at all. The CRediT section credits Dongmei Zhao with "Software" (page 21, line 13), confirming code exists, but no repository URL, no supplementary materials link, and no open-source license are mentioned anywhere in the 22-page paper. So neither the data NOR the code is publicly available.\\\\This is relevant to the PAT-Scan comparison: PAT-Scan's codebase lives in a git repository with version-controlled training scripts, while the Wei method's implementation is effectively inaccessible to other researchers without author cooperation.}\sethlcolor{hlgreen}\hl{Data will be made available on request.}

\vspace{0.5em}

% ---------- Acknowledgments ----------
\noindent\textbf{Acknowledgments}

\vspace{0.3em}

\noindent
The authors acknowledge the support from Innovative Research
Groups of the National Natural Science Foundation of China
(11821202), National Natural Science Foundation of China
(12002075), Liao Ning Revitalization Talents Program
(XLYC1907119), Program for Changjiang Scholars, Innovative
Research Team in University (PCSIRT) and 111 Project (B14013).

\noindent
We also acknowledge Bo Wang and Jiaqi Hong for their helpful
suggestions on the revised manuscript.

\vspace{0.5em}

% ---------- References [1]-[4] ----------
\noindent\textbf{References}

\vspace{0.3em}

\begin{list}{}{%
  \setlength{\leftmargin}{2em}%
  \setlength{\itemindent}{-2em}%
  \setlength{\itemsep}{0.5em}%
  \setlength{\parsep}{0pt}%
}

\item [1] O. Osanai, M. Ohtsuka, M. Hotta, T. Kitaharai, Y.
Takema, A new method for the visualization and quantification
of internal skin elasticity by ultrasound imaging,
\emph{Skin Res. Technol.} 17 (3) (2011) 270--277,
https://doi.org/10.1111/j.1600-0846.2010.00492.x.

\item [2] Y. Zhao, B. Feng, J. Lee, N. Lu, D.M. Pierce, A
multi-layered computational model for wrinkling of human skin
predicts aging effects, \emph{J. Mech. Behav. Biomed. Mater.}
103 (2020) 103552, https://doi.org/10.1016/j.jmbbm.2019.103552.

\item [3] I.L. Kruglikov, P.E. Scherer, Skin aging as a
mechanical phenomenon: the main weak links, \emph{Nutr. Healthy
Aging} 4 (4) (2018) 291--307,
https://doi.org/10.3233/NHA-170037.

\item [4] G.A. Holzapfel, T.C. Gasser, R.W. Ogden, A New
Constitutive Framework for Arterial Wall Mechanics and a
Comparative Study of Material Models, \emph{J. Elast. Phys.
Sci. Solids} 61 (1) (2000) 1--48,
https://doi.org/10.1023/A:1010835316564.

\end{list}

% [page ends here — continues on page 22]

\end{document}
